\chapter{Conclusion}

This project demonstrates that textbook-grounded RAG can improve small local \gls{llm} accuracy on MedQA, but the improvement is conditional on two factors: the embedding model must be biomedical-grade, and the \gls{llm} must be capable of reasoning over retrieved passages. The six-run ablation grid establishes this clearly. Standard Qwen3-4B cannot use retrieved context regardless of embedding, resulting in a -4\gls{pp} penalty. The Thinking variant paired with Octen-Embedding-0.6B achieves 44\% accuracy with RAG, a +11\gls{pp} gain over the 33\% no-\gls{rag} baseline. 

The remaining 30\gls{pp} gap to \gls{amg-rag} (74.1\%, GPT-4o-mini with dynamic PubMed knowledge graph) reflects three differences: \gls{llm} parameter scale, dynamic vs. static knowledge retrieval, and the inherently harder MedQA benchmark. Closing this gap would require scaling the \gls{llm} beyond the T4 \gls{vram} budget, replacing the static textbook corpus with live PubMed querying, and possibly fine-tuning the embedding model on the MedQA retrieval task directly. 

The classification, clustering, \gls{ner}, and recommender components each perform their intended roles. The \gls{tfidf} + \gls{lsvc} classifier achieves F1 = 0.69 on 19 specialties. BERTopic discovers 55 clinically coherent topic groups that provide finer-grained routing than specialty labels alone. \gls{ner} rewriting is validated empirically: it improves retrieval for biomedical embeddings and correctly identifies its own failure mode for general-purpose ones. \gls{knnb} identifies at least one weak specialty for 74\% of students at K=10. Together these components constitute a functional, evaluated, and reproducible medical study agent.